\documentclass[11pt,compress,t,notes=noshow, xcolor=table]{beamer}
\usepackage[]{graphicx}\usepackage[]{color}
% maxwidth is the original width if it is less than linewidth
% otherwise use linewidth (to make sure the graphics do not exceed the margin)
\makeatletter
\def\maxwidth{ %
  \ifdim\Gin@nat@width>\linewidth
    \linewidth
  \else
    \Gin@nat@width
  \fi
}
\makeatother

\newcommand{\citebutton}[2]{%
\beamergotobutton{\href{#2}{#1}}%
}

\newcommand{\blu}[1]{\textcolor{blue}{#1}}
\newcommand{\org}[1]{\textcolor{orange}{#1}}
\newcommand{\ques}{\textbf{\textcolor{red}{Question:  }}}
\newcommand{\questionssofar}{\begin{frame}\frametitle{Any questions?}\end{frame}}

\newcommand\warning{%
 \makebox[1.4em][c]{%
 \makebox[0pt][c]{\raisebox{.1em}{\scriptsize!}}%
 \makebox[0pt][c]{\color{red}\normalsize$\bigtriangleup$}}}%

\definecolor{fgcolor}{rgb}{0.345, 0.345, 0.345}
\newcommand{\hlnum}[1]{\textcolor[rgb]{0.686,0.059,0.569}{#1}}%
\newcommand{\hlstr}[1]{\textcolor[rgb]{0.192,0.494,0.8}{#1}}%
\newcommand{\hlcom}[1]{\textcolor[rgb]{0.678,0.584,0.686}{\textit{#1}}}%
\newcommand{\hlopt}[1]{\textcolor[rgb]{0,0,0}{#1}}%
\newcommand{\hlstd}[1]{\textcolor[rgb]{0.345,0.345,0.345}{#1}}%
\newcommand{\hlkwa}[1]{\textcolor[rgb]{0.161,0.373,0.58}{\textbf{#1}}}%
\newcommand{\hlkwb}[1]{\textcolor[rgb]{0.69,0.353,0.396}{#1}}%
\newcommand{\hlkwc}[1]{\textcolor[rgb]{0.333,0.667,0.333}{#1}}%
\newcommand{\hlkwd}[1]{\textcolor[rgb]{0.737,0.353,0.396}{\textbf{#1}}}%
\let\hlipl\hlkwb

\usepackage{framed}
\makeatletter
\newenvironment{kframe}{%
 \def\at@end@of@kframe{}%
 \ifinner\ifhmode%
  \def\at@end@of@kframe{\end{minipage}}%
  \begin{minipage}{\columnwidth}%
 \fi\fi%
 \def\FrameCommand##1{\hskip\@totalleftmargin \hskip-\fboxsep
 \colorbox{shadecolor}{##1}\hskip-\fboxsep
     % There is no \\@totalrightmargin, so:
     \hskip-\linewidth \hskip-\@totalleftmargin \hskip\columnwidth}%
 \MakeFramed {\advance\hsize-\width
   \@totalleftmargin\z@ \linewidth\hsize
   \@setminipage}}%
 {\par\unskip\endMakeFramed%
 \at@end@of@kframe}
\makeatother

\definecolor{shadecolor}{rgb}{.97, .97, .97}
\definecolor{messagecolor}{rgb}{0, 0, 0}
\definecolor{warningcolor}{rgb}{1, 0, 1}
\definecolor{errorcolor}{rgb}{1, 0, 0}
\newenvironment{knitrout}{}{} % an empty environment to be redefined in TeX

\usepackage{alltt}
\newcommand{\SweaveOpts}[1]{}  % do not interfere with LaTeX
\newcommand{\SweaveInput}[1]{} % because they are not real TeX commands
\newcommand{\Sexpr}[1]{}       % will only be parsed by R
\newcommand{\xmark}{\ding{55}}%


\usepackage[english]{babel}
\usepackage[utf8]{inputenc}

\usepackage{dsfont}
\usepackage{verbatim}
\usepackage{amsmath}
\usepackage{amsfonts}
\usepackage{amssymb}
\usepackage{bm}
\usepackage{csquotes}
\usepackage{multirow}
\usepackage{longtable}
\usepackage{booktabs}
\usepackage{enumerate}
\usepackage[absolute,overlay]{textpos}
\usepackage{psfrag}
\usepackage{algorithm}
\usepackage{algpseudocode}
\usepackage{eqnarray}
\usepackage{arydshln}
\usepackage{tabularx}
\usepackage{placeins}
\usepackage{tikz}
\usepackage{setspace}
\usepackage{colortbl}
\usepackage{mathtools}
\usepackage{wrapfig}
\usepackage{bm}
\usepackage{amsmath}
\usepackage{pifont}

\usetikzlibrary{shapes.multipart,shapes,arrows,automata,positioning,calc,chains,trees, shadows}
\tikzset{
  %Define standard arrow tip
  >=stealth',
  %Define style for boxes
  punkt/.style={
    rectangle,
    rounded corners,
    draw=black, very thick,
    text width=6.5em,
    minimum height=2em,
    text centered},
  % Define arrow style
  pil/.style={
    ->,
    thick,
    shorten <=2pt,
    shorten >=2pt,}
}

\tikzstyle{vec}=[draw, rectangle, fill = white, minimum width=5mm, minimum height=1cm, inner sep = 2pt]

\usepackage{subfig}

% Defines macros and environments
\usepackage{../../style/lmu-lecture}


\let\code=\texttt
\let\proglang=\textsf

\setkeys{Gin}{width=0.9\textwidth}

\setbeamertemplate{frametitle}{\expandafter\uppercase\expandafter\insertframetitle}

\usepackage{bbm}
% basic latex stuff
\newcommand{\pkg}[1]{{\fontseries{b}\selectfont #1}} %fontstyle for R packages
\newcommand{\lz}{\vspace{0.5cm}} %vertical space
\newcommand{\dlz}{\vspace{1cm}} %double vertical space
\newcommand{\oneliner}[1] % Oneliner for important statements
{\begin{block}{}\begin{center}\begin{Large}#1\end{Large}\end{center}\end{block}}


%new environments
\newenvironment{vbframe}  %frame with breaks and verbatim
{
 \begin{frame}[containsverbatim,allowframebreaks]
}
{
\end{frame}
}

\newenvironment{vframe}  %frame with verbatim without breaks (to avoid numbering one slided frames)
{
 \begin{frame}[containsverbatim]
}
{
\end{frame}
}

\newenvironment{blocki}[1]   % itemize block
{
 \begin{block}{#1}\begin{itemize}
}
{
\end{itemize}\end{block}
}

\newenvironment{fragileframe}[2]{  %fragile frame with framebreaks
\begin{frame}[allowframebreaks, fragile, environment = fragileframe]
\frametitle{#1}
#2}
{\end{frame}}


\newcommand{\myframe}[2]{  %short for frame with framebreaks
\begin{frame}[allowframebreaks]
\frametitle{#1}
#2
\end{frame}}

\newcommand{\remark}[1]{
  \textbf{Remark:} #1
}


\newenvironment{deleteframe}
{
\begingroup
\usebackgroundtemplate{\includegraphics[width=\paperwidth,height=\paperheight]{../style/color/red.png}}
 \begin{frame}
}
{
\end{frame}
\endgroup
}
\newenvironment{simplifyframe}
{
\begingroup
\usebackgroundtemplate{\includegraphics[width=\paperwidth,height=\paperheight]{../style/color/yellow.png}}
 \begin{frame}
}
{
\end{frame}
\endgroup
}\newenvironment{draftframe}
{
\begingroup
\usebackgroundtemplate{\includegraphics[width=\paperwidth,height=\paperheight]{../style/color/green.jpg}}
 \begin{frame}
}
{
\end{frame}
\endgroup
}
% https://tex.stackexchange.com/a/261480: textcolor that works in mathmode
\makeatletter
\renewcommand*{\@textcolor}[3]{%
  \protect\leavevmode
  \begingroup
    \color#1{#2}#3%
  \endgroup
}
\makeatother





\input{../../latex-math/basic-math.tex}
\input{../../latex-math/basic-ml.tex}
\newcommand*\POS[1]{\textsubscript{\texttt{#1}}} % tag with part of speech
\usepackage{qtree} %parse tree

\newcommand{\learninggoals}{
\item Understand the scope of the course
\item Answers to all open question
\item Get an impression of the workload}

\title{Basics}
% \author{}
\institute{\href{https://slds-lmu.github.io/lecture_dl4nlp/}{slds-lmu.github.io/lecture\_dl4nlp}}
\date{}

\begin{document}
\lecturechapter{Introduction and Course Outline}
\lecture{Deep Learning for NLP} % stays constant

% ------------------------------------------------------------------------------

\begin{vbframe}{people}

\vfill

\textbf{Lecturers and tutors:}

	\begin{itemize}
		\item Prof. Dr. Hinrich Schütze (weeks 7 \& 9-12)
		\item Dr. Matthias Aßenmacher (weeks 1-6, \& 8)
		\item M.Sc. Leonie Weißweiler (exercise sessions)
		\item Ingo Ziegler (tutoring session)
		\item Marwin Härttrich (exercise correction)
	\end{itemize}

\textbf{Time:}

	\begin{itemize}
		\item Lecture: Wednesday, 9.00 - 11.00
		\item Exercise: Wednesday, 11.00 - 12.00
		\item Tutoring: Thursday, 14.00 - 16.00
	\end{itemize}

\vfill

\end{vbframe}


% ------------------------------------------------------------------------------

\begin{vbframe}{structure of the course}

\vfill

\textbf{Central building blocks of the lecture}

	\begin{itemize}
		\item Basic concepts (2 weeks)
		\item Transformer in-depth (1 week)
		\item BERT (1 week)
		\item BERTology, BERT's successors and distillation (1 week)
		\item Other encoder models \& T5 (1 week)
		\item GPT series and Prompting (2 weeks)
		\item Generative LLMs \& RLHF (1 week)
		\item Multilinguality (1 week)
		\item Math behind training LLMs (1 week)
		\item Evaluation (1 week)
	\end{itemize}

\vfill

\end{vbframe}

% ------------------------------------------------------------------------------

\begin{vbframe}{structure of the exercise}

\vfill

\textbf{Assignments:}

	\begin{itemize}
		\item Training RNNs (week 2 -- 3)
		\item Transformer (week 4 -- 5)
		\item Fine-tuning BERT \& T5 (week 6 -- 7)
		\item Inner workings \& Decoding (week 8 -- 9)
		\item LLama-2 \& Prompting (week 10 -- 11)
	\end{itemize}

\vfill

\end{vbframe}

% ------------------------------------------------------------------------------

\begin{vbframe}{structure of the exercise}

\vfill

\textbf{Central other building blocks of the exercise:}

	\begin{itemize}
		\item Python essentials (week 1)
		\item PyTorch logic (week 3)
		\item Debugging (week 5)
		\item Huggingface universe (week 7)
		\item tbd (week 9)
		\item tbd (week 11)
	\end{itemize}

\vfill

\end{vbframe}

% ------------------------------------------------------------------------------

\begin{vbframe}{prerequisites}
	
\vfill
	
\textbf{Machine learning basics:}

		\begin{itemize}
				\item A proper understanding of\\
				- linear algebra\\
				- loss functions\\
				- regularization\\
				- classification vs. regression\\
				- backpropagation and gradient descent\\
				- simple neural networks (MLPs)
				\item Great resource: \href{https://slds-lmu.github.io/i2ml/}{\textit{I2ML course}}
		\end{itemize}

\vfill

\end{vbframe}

% ------------------------------------------------------------------------------

\begin{vbframe}{Managing expections}

\vfill

\textbf{What to expect}

	\begin{itemize}
		\item A solid and proper understanding of\\
		- the central concepts and models of contemporary NLP\\
		- the capabilities and limitations of different models
		\item Challenging exercises that deepen your unterstanding, i.e.\\
		- you will have to invest quite some time in trying to solve them\\
		- just checking our solutions won't get you far
		\item An (inter)actively taught course with motivated lecturers, i.e.\\
		- we expect you to actively participate in the lecture and ask questions\\
		- you won't get much out of this course by just enrolling and not showing up
	\end{itemize}


\vfill

\end{vbframe}

% ------------------------------------------------------------------------------

\begin{vbframe}{Managing expections}

\vfill

\textbf{What not to expect}

	\begin{itemize}
		\item A prompt engineering course
		\item Riding the hypetrain
		\item Only the latest state-of-the-art LLMs (you also need to know the relevant basics to understand what is going on)
		\item Low workload and easy exercises (you will only learn this thoroughly be investing some time)
	\end{itemize}

\vfill

\end{vbframe}

% ------------------------------------------------------------------------------

\begin{vbframe}{Managing expections}

\vfill

\textbf{Honesty}

	\begin{itemize}
		\item No intention to overload you with (unnecessary) work
		\item Building up proper skills requires \textbf{time investment}
		\item Coding assignments freshly created, we put a lot of work into creating creative and helpful content for you to acquire necessary skills to succeed in this field
		\item Up-front communication of what is going to be required
	\end{itemize}
\vfill

\end{vbframe}

% ------------------------------------------------------------------------------

\begin{vbframe}{Managing expections}

\vfill

\textbf{Workload}

	\begin{itemize}
		\item This is a 6 ECTS course
		\item 1 ECTS := 30h (i.e. this course = 180h)
		\item 13 weeks $\times$ 3h presence $\approx$ 40h
		\item 140h for preparation, self-study, \textbf{working on assignments} 
		\item 140h / 14 weeks = \textbf{10h / week} (\textit{additionally to in-class time})
	\end{itemize}


\vfill

\end{vbframe}

% ------------------------------------------------------------------------------

\begin{vbframe}{your turn}

\vfill

\centering \Huge{\bf Any Questions?}

\vfill

\end{vbframe}

% ------------------------------------------------------------------------------

\endlecture
\end{document}

% ------------------------------------------------------------------------------

\begin{vbframe}{notation (important!!)}

\vfill

\begin{tabular}{ll}
        $\mathcal{C}$ & A corpus of $M$ documents\\[.5em]
        $\mathbf{C}$ & The confusion matrix of a classification problem\\[.5em]
        $C$ & Size of the context window\\[.5em]
        $E$ & Embedding size // Size of the vector representation\\[.5em]
        $h$ & Number of Attention heads in a Transformer-based architecture\\[.5em]
        $H$ & Dimension of the hidden layer of a neural network\\[.5em]
        $i$ & Running index for observations in a data set, words/tokens in a\\
            & sequence or documents in a corpus\\[.5em]
        $j$ & Running index for covariates or levels of a categorical target variable\\[.5em]
        $k$ & The number of classes of a categorical target variable\\[.5em]
        $\mathcal{L}$ & A data generating process a.k.a. language\\[.5em]
        $\mathcal{M}$ & A (language) model\\[.5em]
\end{tabular}

\vfill

\end{vbframe}

% ------------------------------------------------------------------------------

\begin{vbframe}{notation (important!!)}

\vfill

\begin{tabular}{ll}
        $m$ & The number of words/tokens in a sequence $s$\\[.5em]
        $M$ & The number of documents in a corpus $\mathcal{C}$\\[.5em]
        $n$ & The number of observations in a data set\\[.5em]
        $p$ & The number of features in a tabular data set\\[.5em]
        $s$ & An ordered sequence of $m$ words\\[.5em]
        $\theta$ & The parameters of a model\\[.5em]
        $V = \{ w_1, \hdots, w_N \}$ & The vocabulary of a corpus $\mathcal{C}$ or a language $\mathcal{L}$\\[.5em]
        $|V|$ & The size of a vocabulary $V$\\[.5em]
        $w_{\left[t\right]}$ & The word/token at the $t$-th position in a sequence\\[.5em]
				$\vec w^{({\text{word}})}$ or $\vec v^{({\text{word}})}$ & vector representation of a word\\[.5em]
        $\mathbf{W}$ & The weight matrix of a neural network\\[.5em]
        $x$ & Influential variable(s)\\[.5em]
        $y$ & Target variable(s)\\[.5em]
\end{tabular}

\vfill

\end{vbframe}